%\documentclass[12pt,letterpaper]{article}



%%%
%% Required packages:
%%
\usepackage{etex}
\usepackage{amsmath}
\usepackage{amssymb}
\usepackage{amstext}
\usepackage{amsthm}
\usepackage{fancyhdr,fancybox}
\usepackage{mathrsfs} % need for \mathscr command
\usepackage{multicol}
\usepackage[normalem]{ulem}
%\usepackage{pictex,pictexplus}
% \usepackage{pictexwd}
\usepackage{rotating}
\usepackage{url}
\usepackage{xfrac}
\usepackage{booktabs}
\usepackage{color,soul}
\usepackage{comment}

\usepackage{hyperref}
\hypersetup{
    colorlinks=true,     % false: boxed links; true: colored links
    linkcolor=blue,      % color of internal links
    citecolor=blue,      % color of links to bibliography
    filecolor=blue,      % color of file links
    urlcolor=blue        % color of external links
}


\usepackage{tikz}
\usetikzlibrary{backgrounds,calc}

%%
%% Common formatting commands
%%
\setlength{\parindent}{0in}
\setlength{\textwidth}{7in}
\setlength{\evensidemargin}{-0.25in}
\setlength{\oddsidemargin}{-0.25in}
\setlength{\parskip}{.5\baselineskip}
\setlength{\topmargin}{-0.5in}
\setlength{\textheight}{9in}

%               Problem and Part
\newcounter{problemnumber}
\newcounter{partnumber}
\newcounter{subpartnumber}
\newcommand{\Problem}{%
  \refstepcounter{problemnumber}%
  \setcounter{partnumber}{0}%
  \setcounter{subpartnumber}{0}%
  \item[\makebox{\hfil\textbf{\theproblemnumber.}\hfil}]%
  }
\newcommand{\InProblem}[2][1.6]{%
  \refstepcounter{problemnumber}%
  \setcounter{partnumber}{0}%
  \setcounter{subpartnumber}{0}%
  \textbf{\theproblemnumber.}\ \ \parbox[t]{#1in}{#2}%
}
\newcommand{\InSmallProblem}[1]{\InProblem[1.05]{#1}}
\newcommand{\NoProblem}[1]{%
  \mbox{\hfil\textbf{#1}\hfil}%
  }
\newcommand{\UnProblem}{%
  \refstepcounter{problemnumber}%
  \setcounter{partnumber}{0}%
  \setcounter{subpartnumber}{0}%
  \mbox{\hfil\textbf{\theproblemnumber.}\hfil}%
  }
\newcommand{\InFlexProblem}[1]{%
  \refstepcounter{problemnumber}%
  \setcounter{partnumber}{0}%
  \setcounter{subpartnumber}{0}%
  \textbf{\theproblemnumber.}\ \ #1%
}
%%  Part:
\newcommand{\Part}{%
  \renewcommand{\thepartnumber}{(\alph{partnumber})}%
  \refstepcounter{partnumber}
  \setcounter{subpartnumber}{0}%
  \item[(\alph{partnumber})]%
}
\newcommand{\InPart}[2][1.75]{%
  \renewcommand{\thepartnumber}{(\alph{partnumber})}%
  \refstepcounter{partnumber}%
  \setcounter{subpartnumber}{0}%
  (\alph{partnumber})\ \ \parbox[t]{#1in}{#2}%
}
\newcommand{\UnPart}{%
  \refstepcounter{partnumber}%
  \renewcommand{\thepartnumber}{(\alph{partnumber})}%
  \setcounter{subpartnumber}{0}%
  (\alph{partnumber})%
}
\newcommand{\InLargePart}[1]{\InPart[3.05]{#1}}
\newcommand{\InFullPart}[1]{\InPart[6.2]{#1}}
\newcommand{\InSmallPart}[1]{\InPart[1.05]{#1}}
%% SubPart:
\newcommand{\SubPart}{%
  \refstepcounter{subpartnumber}%
  \item[(\roman{subpartnumber})]%
  }
\newcommand{\InSubPart}[2][2.25]{%
  \stepcounter{subpartnumber}%
  (\roman{subpartnumber})\ \ \parbox[t]{#1in}{#2}%
}
\newcommand{\InSmallSubPart}[1]{\InSubPart[1.5]{#1}}
\newcommand{\InTinySubPart}[1]{%
  \stepcounter{subpartnumber}%
  (\roman{subpartnumber})\ \ \makebox{#1}%
}
\newcommand{\InMySubPart}[2][2.25]{%
  \stepcounter{subpartnumber}%
  \parbox[t]{#1in}{\makebox[0.3in][l]{(\roman{subpartnumber})}\ \ {#2}}%
}
\newcommand{\TrueFalse}{\textbf{T}\ \ \textbf{F}\ \ }
\newcommand{\TFPart}[1]{% 
  \stepcounter{partnumber}%
  \setcounter{subpartnumber}{0}%
  \item[(\alph{partnumber})] \TrueFalse\parbox[t]{5.5in}{#1}}

\newcommand{\Points}[1]{(#1 points) \ }
\newcommand{\Point}[1]{(#1 point) \ }


%% For Answers:
\newcounter{answernumber}
\newcommand{\Answer}{%
  \refstepcounter{answernumber}%
  \setcounter{partnumber}{0}%
  \setcounter{subpartnumber}{0}%
  \item[\makebox{\hfil\textbf{\theanswernumber.}\hfil}]%
  }


% Setting up the headers for the first page
%\chead{} % will default to what is typed in the file.
\fancyhead[L]{\textsc{Math 22a}}
\fancyhead[R]{\textsc{Fall 2024}}
\fancyfoot[R]{
	\vspace*{\fill}\footnotesize\textcopyright{} \emph{Harvard Math 22a}	
}
\renewcommand{\headrulewidth}{0pt}

%\fancyhead[C]{}
%	\chead{}	
%	\lhead{}
%	\rhead{}
%	\cfoot{}


% Putting copyright on the footer of each page
%\fancypagestyle{secondpage}{
%	\fancyhead[C]{TESTing again} % change this to be blank
%		%\chead{}	
%	\fancyhead[L]{bears} % change this to be blank
%	\fancyhead[R]{dogs} % change this to be blank
%	\fancyfoot[R]{ %keeps the footer the same
%		\vspace*{\fill}\footnotesize\textcopyright{} \emph{Harvard Math 22a}	
%	}
%}
%\renewcommand{\headrulewidth}{0pt}
%\pagestyle{fancy}
\pagestyle{empty}


%\lhead{}
%\rhead{}
%\rfoot{\vspace*{\fill}\footnotesize\textcopyright{} \emph{Harvard Math 22a}}
%\renewcommand{\headrulewidth}{0pt}
%\pagestyle{fancy}




%%
%% Common shortcut commands:
%%
\newcommand{\ds}{\displaystyle}
\newcommand{\sgn}{\operatorname{sgn}}
\renewcommand{\Pr}{\operatorname{P}}
\newcommand{\CPr}[3][{}]{\Pr_{\text{#1}}\left(#2\,|\,#3\right)}

\newcommand{\C}{\mathbb{C}}
\newcommand{\R}{\mathbb{R}}
\newcommand{\Z}{\mathbb{Z}}
\newcommand{\N}{\mathbb{N}}
\newcommand{\Q}{\mathbb{Q}}
\newcommand{\D}{\mathbb{D}}

\newcommand{\rref}{\operatorname{rref}}
\newcommand{\im}{\operatorname{im}}
\newcommand{\rank}{\operatorname{rank}}
\newcommand{\diag}{\operatorname{diag}}
\newcommand{\Span}{\operatorname{span}}
%\renewcommand{\vec}[1]{\mathbf{#1}}
\newcommand{\charpoly}{\mathscr{P}}
\newcommand{\nicefrac}[2]{\sfrac{#1}{#2}} % using xfrac package
\newcommand{\topology}{\mathcal{T}}
\newcommand{\Inn}{\operatorname{Inn}}

\newcommand{\blankbox}[2]{\fbox{\rule{#1}{0in}\rule{0in}{#2}}}

\newenvironment{myindentpar}[1]%
 {\begin{list}{}%
         {\setlength{\leftmargin}{#1}
          \setlength{\rightmargin}{\leftmargin}}%
         \item[]%
 }
 {\end{list}}
\newtheorem{theorem}{Theorem}
\newtheorem*{NStheorem}{Theorem}
\newtheorem{cor}[theorem]{Corollary}
\newtheorem*{claim}{Claim}
\newtheorem{defn}{Definition}

\newenvironment{amatrix}[1]{%
  \left(\begin{array}{@{}*{#1}{c}|c@{}}
}{%
  \end{array}\right)
}

%--------grstep
% For denoting a Gauss' reduction step.
% Use as: \grstep{\rho_1+\rho_3} or \grstep[2\rho_5 \\ 3\rho_6]{\rho_1+\rho_3}
\newcommand{\grstep}[2][\relax]{%
   \ensuremath{\mathrel{
       {\mathop{\longrightarrow}\limits^{#2\mathstrut}_{
                                     \begin{subarray}{l} #1 \end{subarray}}}}}}
\newcommand{\swap}{\leftrightarrow}




\section{{Matrix Operations, $\vec{\S}$2.1, 2.2}}% 
\chead{\Large\textbf{Matrix Operations, $\vec{\S}$2.1, 2.2}}% 

%\begin{document}
\thispagestyle{fancy}

Recall from last class:

\begin{itemize}

\item {\bf Theorem 10:} Let $T: \R^n \to \R^m$ be linear. Then there exists a unique matrix $A$ such that $T(\vec{x})=A\vec{x}$ for all $\vec{x} \in \R^n$; namely $A=[T(\vec{e}_1) \; \dots \; T(\vec{e}_n)]$.
\item {\bf Theorem 11:} Let $T: \R^n \to \R^m$ be linear. Then $T$ is injective if and only if $T(\vec{x})=\vec{0}$ has only the trivial solution.
\item {\bf Theorem 12:} Let $T: \R^n \to \R^m$ be linear, and let $A$ be the associated matrix. Then 
\begin{enumerate}
\item $T$ is surjective if and only if the columns of $A$ span $\R^m$.
\item $T$ is injective if and only if the columns of $A$ are linearly independent. 
\end{enumerate}
\end{itemize}

\newpage

Let $A=[\vec{a}_1 \; \dots \; \vec{a}_n]$ be an $m$ by $n$ matrix. We express $A$ by
$$A=\begin{bmatrix}
a_{11} & a_{12} & \hdots & a_{1n}\\ 
a_{21} & a_{22} & \hdots & a_{2n}\\ 
\vdots & & \ddots & \vdots\\
a_{n1} & a_{n2} & \hdots & a_{nn}\\ 
\end{bmatrix}$$
where $a_{ij}$ is the entry in the $i$-th row and $j$-th column.

{\bf Theorem:} Let $A$, $B$, and $C$ be matrices of the same size, and let $r$ and $s$ be real numbers. Then the following properties hold:
\begin{enumerate}
\item $A+B=B+A$.
\item $(A+B)+C=A+(B+C)$.
\item $A+0=A$.
\item $r(A+B)=rA+rB$.
\item $(r+s)A=rA+sA$.
\item $r(sA)=rsA$.
\end{enumerate}

\begin{defn}
If $A$ is a $m\times n$ matrix and $B$ is an $n\times p$ matrix with columns $\vec{b}_1, \dots, \vec{b}_p$, then $AB$ is the $m \times p$ matrix given by $$AB=A[\vec{b}_1 \; \dots \; \vec{b}_p]=[A\vec{b}_1 \; \dots \; A\vec{b}_p].$$
\end{defn}

\begin{enumerate}

\Problem Compute $AB$ when $$A=\begin{bmatrix} 2 & 3\\ 1 & -5\\ \end{bmatrix} \text{ and } B=\begin{bmatrix} 4 & 3 & 6\\ 1 & -2 & 3\\ \end{bmatrix}.$$ Can you also compute $BA$?

\newpage


{\bf Remark:}
$$(AB)_{ij}=a_{i1}b_{1j}+a_{i2}b_{2j}+\dots+a_{in}b_{nj}.$$

{\bf Properties of Matrix Multiplication.}

{\bf Theorem:}
\begin{enumerate}
\item $A(BC)=(AB)C$.
\item $A(B+C)=AB+AC$.
\item $(B+C)A=BA+CA$.
\item $r(AB)=(rA)B=A(rB)$.
\item $I_m A =A = AI_n$.
\end{enumerate}

\Problem {\bf True or False Questions.} Prove or give a counterexample as appropriate. Suppose $A$ and $B$ are $n$ by $n$ matrices.

\begin{enumerate}
\item $AB=BA$.

\vfill

\item If $AB=0$, then $A=0$ or $B=0$.

\vfill

\item If $A^2=0$, then $A=0$.

\vfill

\item If $AB=AC$, then $B=C$.

\vfill

\end{enumerate}

\newpage


\begin{defn}
Given an $m\times n$ matrix $A$, then the {\bf transpose of $A$}, denoted $A^{T}$ is the $n \times m$ matrix given by taking the columns of $A$ as rows.
\end{defn}

{\bf Theorem:} Let $A$ be an $m$ by $n$ matrix, $B$ be an $n$ by $p$ matrix, and $r\in \R$.
\begin{enumerate}
\item $(A^{T})^{T}=A.$
\item $(A+B)^{T}=A^{T}+B^{T}.$
\item $(rA)^{T}=r A^{T}$.
\item $(AB)^{T}=B^{T} A^{T}$.
\end{enumerate}

\begin{defn}
An $n$ by $n$ matrix $A$ is {\bf invertible} if there is an $n$ by $n$ matrix $C$ so that $AC=I$ and $CA=I$. We denote the inverse of $A$ by $A^{-1}$.
\end{defn}

{\bf Natural Questions:}
\begin{enumerate}
\item Do inverses always exist?
\item If not, when do they exist?
\item If they exist, are they unique?
\item How can we find them?
\end{enumerate}

\Problem Let $A=\begin{bmatrix} a & b\\ c & d\\ \end{bmatrix}$. Determine when $A$ is invertible and find $A^{-1}$. 

\newpage

{\bf Theorem:} Let $A=\begin{bmatrix} a & b \\ c & d\\ \end{bmatrix}$. If $ad-bc\neq 0$, then $A$ is invertible  and $$A^{-1}= \frac{1}{ad-bc} \begin{bmatrix} d & -b\\ -c &a \\ \end{bmatrix}.$$

\begin{defn}
We define the {\bf determinant} of $A=\begin{bmatrix} a & b \\ c & d\\ \end{bmatrix}$, denoted by $\text{det}(A)$, by $$\text{det}(A)=ad-bc.$$
\end{defn}

{\bf Theorem:} If $A$ is an $n$ by $n$ invertible matrix, then for each $\vec{b}\in \R^n$, $A\vec{x}=\vec{b}$ has a unique solution; namely $\vec{x}=A^{-1}\vec{b}$.

\Problem Prove the Theorem.

\vfill

\newpage

{\bf Theorem:} Suppose $A$ and $B$ are $n$ by $n$ matrices, then the following hold. 
\begin{enumerate}
\item If $A$ is invertible, then $A^{-1}$ is invertible and $(A^{-1})^{-1}=A$.
\item $(AB)^{-1}=B^{-1} A^{-1}$.
\item $(A^{T})^{-1}=(A^{-1})^{T}$.
\end{enumerate}

\Problem Prove the Theorem.

 

\end{enumerate}

\begin{comment}

\newpage
\chead{\Large\textbf{Matrix Operations -- Answers / Solutions}}%
\lhead{}
\rhead{}
\thispagestyle{fancy}

\begin{enumerate}
\Answer Let $$A=\begin{bmatrix} 2 & 3\\ 1 & -5\\ \end{bmatrix} \text{ and } B=\begin{bmatrix} 4 & 3 & 6\\ 1 & -2 & 3\\ \end{bmatrix},$$ then using the definition of matrix multiplication we get the following
$$A\vec{b}_1 =\begin{bmatrix}
2(4)+3(1)\\
1(4)-5(1)\\
\end{bmatrix}=\begin{bmatrix}
11\\
-1\\
\end{bmatrix}, A\vec{b}_2 =\begin{bmatrix}
2(3)+3(-2)\\
1(3)-5(-2)\\
\end{bmatrix}=\begin{bmatrix}
0\\
13\\
\end{bmatrix}, \text{ and } A\vec{b}_3 =\begin{bmatrix}
2(6)+3(3)\\
1(6)-5(3)\\
\end{bmatrix}=\begin{bmatrix}
21\\
-9\\
\end{bmatrix}.$$

Hence $$AB=[A\vec{b_1} \; A\vec{b}_2 \; A\vec{b}_3]=\begin{bmatrix}
11 & 0 & 21\\
-1 & 13 & -9\\
\end{bmatrix}.$$

We cannot multiply the matrices in the other order as the number of rows of $A$ does not equal the number of columns of $B$, and hence $BA$. In terms of the corresponding linear transformations, this says the number of outputs of $A$ does not match the number of inputs of $B$, and hence we cannot compose the functions in that order.
  
\Answer   

\begin{enumerate}

\item {\bf False.} There are many counterexamples, but take for instance $A=\begin{bmatrix} 1 & 2\\ 3 & 4\\ \end{bmatrix}$, and $B=\begin{bmatrix} 0 & 1\\ -1 & 1\\ \end{bmatrix}$, then $AB=\begin{bmatrix} -2 & 3\\ -4 & 7\\ \end{bmatrix}$ and $BA=\begin{bmatrix} 3 & 4\\ 2 & 2\\ \end{bmatrix}$. Thus $AB\neq BA$.

\item {\bf False.} Let $A=\begin{bmatrix} 1 & 1\\ 2 & 2\\ \end{bmatrix}$, and $B=\begin{bmatrix} 1 & 1\\ -1 & -1\\ \end{bmatrix}$, then $AB=\begin{bmatrix} 0 & 0\\ 0 & 0\\ \end{bmatrix}$.

\item {\bf False.} Let $A=\begin{bmatrix} 0 & 1\\ 0 & 0\\ \end{bmatrix}$, then $A^2=\begin{bmatrix} 0 & 0\\ 0 & 0\\ \end{bmatrix}$.

\item {\bf False.} Let $$A=\begin{bmatrix} 0 & 1\\ 0 & 1\\ \end{bmatrix}, B=\begin{bmatrix} 1 & 1\\ 0 & 1\\ \end{bmatrix}\text{, and  } C=\begin{bmatrix} 2 & 1\\ 0 & 1\\ \end{bmatrix},$$  then $AB=AC$ but $B\neq C$.

\end{enumerate}
    
\Answer We want to find $e$, $f$, $g$ and $h$ such that 
$$\begin{bmatrix}
a & b\\
c & d\\
\end{bmatrix}\begin{bmatrix}
e & f\\
g &h\\
\end{bmatrix}=\begin{bmatrix}
1 & 0\\
0 & 1\\
\end{bmatrix}.$$
Thus we get the following system of equations:
\begin{align*}
ae+gb&=1\\
af+bh&=0\\
ce+dg&=0\\
cf+dh&=1\\
\end{align*}

Multplying the first equation by $c$ and the third equation by $-a$ and adding them together yields: $bcg-adg=c$, and assuming $ad-bc\neq0$, we get $g=\frac{-c}{ad-bc}$.

Continuing in a similar manner we get $e=\frac{d}{ad-bc}$, $h=\frac{a}{ad-bc}$, and $f=\frac{-b}{ad-bc}$.

Thus the matrix is invertible if and only if $ad-bc\neq0$, and we have an explicit formula for the inverse matrix.
  
\Answer Since $A$ is invertible, we know $A^{-1}$ exists. Let $\vec{x}=A^{-1}\vec{b}$, then we check $\vec{x}$ is a solution. Here $A\vec{x}=AA^{-1}\vec{b}=\vec{b}$, and hence $\vec{x}$ is a solution.

Now we check it is the unique solution. Suppose $\vec{u}$ is another solution, then $A\vec{u}=\vec{b}$. Since $A$ is invertible, we can multiply both sides by $A^{-1}$ to get $A^{-1}A\vec{u}=A^{-1}\vec{b}$, and hence $\vec{u}=A^{-1}\vec{b}=\vec{x}$. Thus $\vec{x}$ is the only solution.
  
\Answer 

\begin{enumerate}
	
\Part Suppose $A$ is invertible. We check that $A$ is the inverse of $A^{-1}$; namely, $AA^{-1}=I_n$ and $A^{-1}A=I_n$. Thus the inverse of $A^{-1}$ is $A$; in symbols $(A^{-1})^{-1}=A$.

\Part Again suppose $A$ and $B$ are both invertible matrices. We check that $B^{-1}A^{-1}$ is the inverse of $AB$: $ABB^{-1}A^{-1}=I_n$ and $B^{-1}A^{-1}AB=I_n$. Thus by the definition of inverse matrix, $B^{-1}A^{-1}$ is the inverse of $AB$; i.e. $(AB)^{-1}=B^{-1}A^{-1}$.

\Part Suppose $A$ is an invertible matrix. We want to prove that the inverse of $A^{T}$ is $(A^{-1})^{T}$. Again we check the conditions in the definition of matrix inverse: $A^{T}(A^{-1})^{T}=(A^{-1}A)^{T}=(I_n)^{T}=I_n$ and $(A^{-1})^{T}A^{T}=(AA^{-1})^{T}=(I_n)^{T}=I_n$ where we used $(AB)^{T}=B^{T}A^{T}$. Thus $(A^{T})^{-1}=(A^{-1})^{T}$.

\end{enumerate}
\end{enumerate}

\end{comment}
%\end{document}



%%% Local ables: 
%%% mode: latex
%%% TeX-master: t
%%% End: 
